\documentclass[a4paper,11pt]{article}
\usepackage{lmodern}
\usepackage[T1]{fontenc}
\usepackage[utf8]{inputenc}
\usepackage{amsmath,amssymb,amsthm}
\usepackage{mathtools}
\usepackage{booktabs}
\usepackage{longtable}
\usepackage{array}
\usepackage{hyperref}
\usepackage{graphicx}
\usepackage{float}
\usepackage{geometry}
\geometry{margin=25mm}

\hypersetup{
  colorlinks=true,
  linkcolor=blue,
  urlcolor=blue,
  citecolor=blue
}

\theoremstyle{definition}
\newtheorem{theorem}{Theorem}[section]
\newtheorem{proposition}[theorem]{Proposition}
\newtheorem{lemma}[theorem]{Lemma}
\newtheorem{corollary}[theorem]{Corollary}
\newtheorem{definition}[theorem]{Definition}
\newtheorem{remark}[theorem]{Remark}
\newtheorem{observation}[theorem]{Observation}

\setlength{\parindent}{1em}
\setlength{\parskip}{0.5em}

\providecommand{\tightlist}{\setlength{\itemsep}{0pt}\setlength{\parskip}{0pt}}
\providecommand{\real}[1]{#1}
\begin{document}

\section{Synthesis and the 4+1 to 3+1 Reduction as the Principal Open
Problem (Paper 6:
Synthesis)}\label{synthesis-and-the-41-to-31-reduction-as-the-principal-open-problem-paper-6-synthesis}

\textbf{Author}: Noriaki Kihara (WF System Co., Ltd.) \textbf{ORCID}:
\href{https://orcid.org/0009-0004-6753-4020}{0009-0004-6753-4020}
\textbf{Date}: April 2026 \textbf{DOI}:
\href{https://doi.org/10.5281/zenodo.19839400}{10.5281/zenodo.19839400}

\begin{center}\rule{0.5\linewidth}{0.5pt}\end{center}

\subsection{Abstract}\label{abstract}

We synthesize the results of the five companion papers and identify the
principal open problem of the program: how the 4+1-dimensional black
hole thermodynamics established by the program (with lifetime
\(\tau \propto M^{3/2}\)) reduces to the observed 3+1-dimensional
thermodynamics (with \(\tau \propto M^3\)). We discuss three candidate
mechanisms --- Kaluza--Klein compactification, brane localization à la
Randall--Sundrum/ADD, and a holographic-style projection --- and propose
a programme of further investigation. The paper contains no new
theorems; its purpose is to crystallize what has and has not been
established, and to point toward the next stage.

\begin{center}\rule{0.5\linewidth}{0.5pt}\end{center}

\subsection{§1. Introduction}\label{introduction}

The program presented in Papers 1 through 5 establishes a
self-consistent geometric and combinatorial framework for black hole
thermodynamics in 4+1 dimensions. From the two hypotheses of Paper 1 ---
the projection hypothesis and the discrete drift hypothesis --- together
with the geometric construction of Paper 2 and the combinatorial result
of Paper 3, the program derives:

\begin{itemize}
\tightlist
\item
  A Schwarzschild--de Sitter-form metric on a four-dimensional
  subjective chart.
\item
  An asymptotic constant \(c = 8/(3\pi)\) for the volume deficit
  \(\Delta(R)\).
\item
  A correspondence with the 4+1-dimensional Tangherlini
  Bekenstein--Hawking entropy with ratio
  \(\Delta(R)/S_{BH} \to 32/(3\pi)\).
\item
  The Tangherlini scaling laws \(R \propto M^{1/2}\),
  \(T_H \propto M^{-1/2}\), \(\tau \propto M^{3/2}\),
  \(dM/dt \propto -M^{-1/2}\).
\end{itemize}

The successes are concentrated in 4+1 dimensions. The principal open
problem, identified in Paper 5 §5, is the relation between this 4+1
structure and the 3+1 physics of observable black holes. We devote this
final paper to that problem.

\begin{center}\rule{0.5\linewidth}{0.5pt}\end{center}

\subsection{§2. Synthesis: From Hypotheses to
Results}\label{synthesis-from-hypotheses-to-results}

The two hypotheses of Paper 1 are partially supported by the companion
papers:

\textbf{Hypothesis 1 (Projection Hypothesis):} A black hole in the
four-dimensional theory is the projected image of a higher-dimensional
structure.

\begin{itemize}
\tightlist
\item
  Paper 2 constructs the gnomonic projection from
  \(S^4(R) \subset \mathbb{R}^5\) and derives the induced
  four-dimensional metric. The metric satisfies
  \(G_{\mu\nu} + (3/R^2) g_{\mu\nu} = 0\) and, in the Schwarzschild--de
  Sitter extension with a mass parameter, exhibits horizon structure.
  The projection thus has, at the level of the metric, the right form to
  support the hypothesis.
\item
  Paper 4 shows that the geometric quantity \(\Delta(R)\) (a property of
  the projection) has the right \(R^3\) scaling to match the 4+1
  Bekenstein--Hawking entropy.
\item
  Paper 5 shows that the curvature radius \(R\) --- the geometric scale
  of the projection --- is determined self-consistently from the matter
  content via \(E_{\mathrm{total}} \propto R^2\), giving
  \(R \propto M^{1/2}\) as in 4+1 Tangherlini.
\end{itemize}

\textbf{Hypothesis 2 (Discrete Drift Hypothesis):} Hawking evaporation
corresponds to the drift of discrete units across a higher-dimensional
region's boundary.

\begin{itemize}
\tightlist
\item
  Paper 3 establishes the combinatorial framework --- unit-cube packing
  in a four-dimensional ball --- and computes the asymptotic volume
  deficit \(\Delta(R) \sim (16\pi/3) R^3\) with leading constant
  \(c = 8/(3\pi)\).
\item
  Paper 4 identifies \(\Delta(R)\) as the candidate ``discrete drift
  entropy.''
\item
  Paper 5 derives the evaporation dynamics from the scaling of \(R\)
  with \(M\), recovering the Tangherlini lifetime
  \(\tau \propto M^{3/2}\).
\end{itemize}

The hypotheses are supported at the level of scaling and structural
form. They are \textbf{not} supported at the level of:

\begin{itemize}
\tightlist
\item
  An identification of \(\Delta(R)\) with a microstate count (the
  constant ratio \(32/(3\pi)\) remains unexplained at the microscopic
  level).
\item
  A complete derivation of the discrete drift mechanism (the program
  assumes that discrete units drift across boundaries; it does not
  derive this from first principles).
\item
  A reduction to physical 3+1-dimensional black hole physics (the
  principal open problem).
\end{itemize}

\begin{center}\rule{0.5\linewidth}{0.5pt}\end{center}

\subsection{§3. The Mathematical--Physical
Bridge}\label{the-mathematicalphysical-bridge}

The program operates in two registers. Papers 2 and 3 are mathematical
--- theorems and proofs about the central projection and the
lattice-packing problem. Papers 4 and 5 are physical --- they take the
mathematical structures and ask whether they correspond to features of
black hole thermodynamics.

The bridge between the two registers is the identification of the
geometric quantity \(R\) with the physical quantity \(r_h\). This
identification is not a mathematical theorem; it is a hypothesis. It is
supported by the agreement of the resulting scaling laws with
Tangherlini, which is a non-trivial check, but it is not derived from
first principles.

The strength of the bridge is therefore the strength of the hypothesis.
If one accepts the hypothesis, the entire framework is internally
consistent and reproduces a substantial body of established results in
4+1-dimensional black hole physics. If one rejects the hypothesis, the
mathematical content of Papers 2 and 3 still stands, but its physical
interpretation is open.

\begin{center}\rule{0.5\linewidth}{0.5pt}\end{center}

\subsection{§4. The Principal Open Problem: Reduction to 3+1
Dimensions}\label{the-principal-open-problem-reduction-to-31-dimensions}

\subsubsection{§4.1 Statement of the
problem}\label{statement-of-the-problem}

The 4+1-dimensional Tangherlini results recovered by the program are:

\begin{itemize}
\tightlist
\item
  \(r_h \propto M^{1/2}\)
\item
  \(T_H \propto M^{-1/2}\)
\item
  \(\tau \propto M^{3/2}\)
\end{itemize}

The 3+1-dimensional Schwarzschild results are:

\begin{itemize}
\tightlist
\item
  \(r_s \propto M\)
\item
  \(T_H \propto M^{-1}\)
\item
  \(\tau \propto M^3\)
\end{itemize}

These are not quantitatively reconcilable without an additional
mechanism. The program, as it stands, describes 4+1-dimensional black
holes; it does not describe 3+1-dimensional black holes.

If the program is to make contact with observable physics, this gap must
be bridged. The problem is precisely formulated:

\begin{quote}
\textbf{Open Problem.} Identify a mechanism by which the 4+1-dimensional
thermodynamics of the central-projection construction reduces to the
3+1-dimensional thermodynamics of physical black holes, with
quantitative agreement at the level of the mass--temperature relation,
the entropy formula, and the evaporation lifetime.
\end{quote}

\subsubsection{§4.2 Candidate mechanism A: Kaluza--Klein
compactification}\label{candidate-mechanism-a-kaluzaklein-compactification}

If one of the four spatial dimensions is compact at a scale \(\ell\)
much smaller than \(r_h\), the black hole horizon wraps around the
compact direction, and the effective 3+1-dimensional theory has
thermodynamics that may differ from the 4+1 case.

In Kaluza--Klein theory, the 4+1-dimensional Schwarzschild--Tangherlini
black hole reduces, for \(r_h \ll \ell\) (``small black holes''), to a
4+1-dimensional black string wrapping the compact direction; for
\(r_h \gg \ell\) (``large black holes''), it reduces to a
3+1-dimensional Schwarzschild black hole with effective Newton constant
\(G_4 = G_5/\ell\).

In the latter regime, the 3+1 thermodynamics is recovered. The
transition between regimes occurs at \(r_h \sim \ell\), the
compactification scale.

For the central-projection construction to produce 3+1 thermodynamics
for astrophysical black holes, \(\ell\) must be larger than the
Schwarzschild radius of those black holes --- which would require
\(\ell > r_s \sim\) km, easily ruled out by observational bounds.

Conversely, \(\ell\) small (Planck scale) would give 4+1 thermodynamics
for all macroscopic black holes, contradicting the 3+1 observed
behaviour.

\textbf{This mechanism does not, by itself, resolve the discrepancy.}

\subsubsection{§4.3 Candidate mechanism B: Brane localization
(Randall--Sundrum /
ADD)}\label{candidate-mechanism-b-brane-localization-randallsundrum-add}

In Randall--Sundrum-type models, our 3+1-dimensional spacetime is a
brane embedded in a higher-dimensional bulk. Matter is localized on the
brane while gravity propagates in the bulk. Black holes on the brane
have a 3+1-dimensional event horizon but bulk gravitational effects.

Such models can produce a hierarchy where: - Large (astrophysical) black
holes are well-approximated by 3+1 Schwarzschild. - Small black holes
(\(r_h \lesssim\) AdS radius) deviate, becoming bulk-dominated and
showing 4+1-like thermodynamics.

This mechanism could in principle accommodate the central-projection
construction: the 4+1 results would apply to small (Planck-scale) black
holes, while the 3+1 results recover for large black holes. The
transition scale would be set by the AdS or compactification radius.

For this to be a viable mechanism in the central-projection program, the
relation between the curvature radius \(R\) of the projection and the
compactification scale would need to be made precise. \textbf{後で検討}.

\subsubsection{§4.4 Candidate mechanism C: Holographic-style
projection}\label{candidate-mechanism-c-holographic-style-projection}

A more speculative mechanism: the 4+1 central-projection structure is a
``bulk'' description, and the observable 3+1 physics is its boundary
projection in the holographic sense. In this picture:

\begin{itemize}
\tightlist
\item
  4+1 dimensions: \((x, y, z, t, R)\) with the Tangherlini-like
  dynamics.
\item
  3+1 dimensions: the boundary of the central-projection structure at
  fixed \(R\) (or asymptotic \(R\)), with thermodynamics derived from
  the 4+1 bulk via a holographic dictionary.
\end{itemize}

The challenge is to specify the holographic dictionary. Standard AdS/CFT
does not directly apply because the central-projection metric is not
asymptotically AdS; it is a Schwarzschild-de Sitter form with
cosmological horizon. A de Sitter / dS-CFT holography would be the
analogous framework, but dS-CFT is itself less developed than AdS/CFT
and lacks a fully accepted formulation.

This mechanism would also place the central-projection program in
dialogue with the holographic principle, which the program initially
declined to invoke. The dialogue would be on equal terms: not asserting
holography, but exploring how the central-projection picture might be
related to it.

\subsubsection{§4.5 What the open problem
requires}\label{what-the-open-problem-requires}

To resolve the open problem requires:

\begin{enumerate}
\def\labelenumi{\arabic{enumi}.}
\tightlist
\item
  A quantitative specification of the mechanism (Kaluza--Klein, brane
  localization, or holographic projection --- or some other).
\item
  Derivation of the 3+1 mass--temperature, mass--lifetime, and
  entropy--area relations from this mechanism applied to the
  central-projection construction.
\item
  Comparison with the standard 3+1 Schwarzschild--Hawking results, with
  quantitative agreement (or with explicit identification of any
  residual discrepancies).
\end{enumerate}

None of these are accomplished by the current program. The problem is
identified, the candidate mechanisms are listed, and the further
investigation is left for future work.

\begin{center}\rule{0.5\linewidth}{0.5pt}\end{center}

\subsection{§5. Conclusion}\label{conclusion}

The program in five companion papers establishes a self-contained
framework for 4+1-dimensional black hole thermodynamics from minimal
geometric and combinatorial assumptions. The successes are real: the
asymptotic constant \(c = 8/(3\pi)\), the entropy ratio
\(\Delta(R)/S_{BH} \to 32/(3\pi)\), and the Tangherlini scaling laws are
all derived without invoking holography, AdS/CFT, or any specific
quantum gravity programme.

The principal open problem is the reduction from 4+1 to 3+1 dimensions.
The program does not solve it; it identifies it precisely and offers
candidate mechanisms whose detailed development is left for future work.

We close with three observations.

First, the program has the methodological virtue that its open problems
are sharp. Where the program agrees with established results, it agrees
quantitatively. Where it does not agree, the disagreement is in a
specific quantitative form (the 4+1 vs 3+1 thermodynamics) that future
work can attempt to address. There is no vague gesture toward ``more
research needed''; there is a specific question.

Second, the program does not require any additional input from quantum
field theory, quantum gravity, or the standard Hawking computation. Its
4+1 thermodynamics is derived from the geometric construction of Paper 2
and the combinatorial result of Paper 3, with the Bekenstein--Hawking
entropy formula as the only external input (and only in Paper 4 for
comparison purposes).

Third, the program's value, regardless of whether the open problem is
eventually resolved, is in the precision of the mathematical content.
Papers 2 and 3 are independent mathematical results; they can be
evaluated on their own terms. Paper 3, in particular, contributes a new
asymptotic determination --- \(c = 8/(3\pi)\) for the volume deficit ---
that has interest within number theory and discrete geometry,
independent of any black hole interpretation.

The program, as presented in these six papers, is complete as a unit.
Whether it leads to further developments --- including a resolution of
the 4+1 to 3+1 problem --- is a question for future work.

\begin{center}\rule{0.5\linewidth}{0.5pt}\end{center}

\subsection{References}\label{references}

\begin{enumerate}
\def\labelenumi{\arabic{enumi}.}
\tightlist
\item
  T. Kaluza, ``Zum Unitätsproblem der Physik,'' \emph{Sitzungsber.
  Preuss. Akad. Wiss.}, 966--972 (1921).
\item
  O. Klein, ``Quantentheorie und fünfdimensionale Relativitätstheorie,''
  \emph{Z. Phys.} \textbf{37}, 895 (1926).
\item
  F. R. Tangherlini, ``Schwarzschild Field in \(n\) Dimensions and the
  Dimensionality of Space Problem,'' \emph{Nuovo Cimento} \textbf{27},
  636 (1963).
\item
  R. C. Myers, M. J. Perry, ``Black Holes in Higher Dimensional
  Spacetimes,'' \emph{Annals Phys.} \textbf{172}, 304 (1986).
\item
  J. Maldacena, ``The Large-\(N\) Limit of Superconformal Field Theories
  and Supergravity,'' \emph{Adv. Theor. Math. Phys.} \textbf{2}, 231
  (1998).
\item
  N. Arkani-Hamed, S. Dimopoulos, G. Dvali, ``The Hierarchy Problem and
  New Dimensions at a Millimeter,'' \emph{Phys. Lett. B} \textbf{429},
  263 (1998).
\item
  L. Randall, R. Sundrum, ``A Large Mass Hierarchy from a Small Extra
  Dimension,'' \emph{Phys. Rev.~Lett.} \textbf{83}, 3370 (1999).
\item
  A. Chamblin, S. W. Hawking, H. S. Reall, ``Brane-world Black Holes,''
  \emph{Phys. Rev.~D} \textbf{61}, 065007 (2000).
\item
  A. Strominger, ``The dS/CFT Correspondence,'' \emph{JHEP}
  \textbf{0110}, 034 (2001).
\item
  R. Emparan, H. S. Reall, ``Black Holes in Higher Dimensions,''
  \emph{Living Rev.~Rel.} \textbf{11}, 6 (2008).
\end{enumerate}

\end{document}
